\chapter{Faget som ikkje kan seie nei}

\begin{motto}
Det gjeldande rammeverket har inga grense for korleis ein kan handsame eit objekt,\\ uavhengig av kva objektet er. Òg når objektet er eit menneske.
\end{motto}

Skaden i førre kapittel kunne ein lese som ein eigenskap ved teknologibransjen, ikkje ved designfaget. Dette kapitlet stengjer den utvegen. For mangelen på fundament er ikkje berre ein epistemisk pinleg detalj; han har ein moralsk konsekvens, og konsekvensen er at faget strukturelt ikkje kan seie nei til noko. Ikkje fordi designarar er dårlege menneske, men fordi eit fag utan ein presis modell av kva ei form gjer — kva agentar ho rører, kva ho ber, kva ho legg til side — manglar den interne målestokken som ville late det skilje forma som ber verda med seg frå forma som sviktar langs alt han ikkje såg. Det som ikkje kan målast, kan ikkje avvisast.

\section{Når objektet er eit menneske}

Tenestedesign lagar og optimaliserer system for målretting, profilering og deportasjon, og kallar det funksjonelt. Palantir er eit tenestedesignprosjekt. Kunden er staten; «brukaren», i den grad ordet tyder noko her, er den som vert overvaka, skåra, sortert, fjerna. Systemet vart bygt for etterretning, selt til politi og grensekontroll, og brukt til systematisk registrering og forfølging av innvandrarfamiliar, ned til born \autocite{biddle_intercept}. Variantar av same arkitektur styrer målval i militære operasjonar; kvart mål er eit punkt i eit formrom optimert for éin akse \autocite{abraham2024lavender}. Systema er funksjonelle. Dei er presist funksjonelle langs den eine aksen dei vart bedne om å optimalisere.

Dette er ikkje ytterkanten av faget; det er sentrum av det, gjort synleg. Virginia Eubanks har dokumentert korleis automatiserte system profilerer, politiovervaker og straffar dei fattige \autocite{eubanks2018}; Ruha Benjamin korleis teknisk «nøytralitet» kodar inn rasehierarki \autocite{benjamin2019}; Cathy O'Neil korleis udokumenterte skåringsmodellar formar liv i stor skala \autocite{oneil2016}. I kvart tilfelle er det snakk om grensesnitt, flytdiagram, datamodellar og avgjerdssystem — det vil seie design. Det gjeldande rammeverket i faget har inga grense for korleis ein kan handsame eit objekt, uavhengig av kva objektet er. Så lenge handsaminga kan kallast funksjonell, er ho innanfor — òg når objektet er levande, òg når objektet er eit menneske. Eit fag som ikkje kan seie kva ein agent er, kva eit seleksjonstrykk er, kva som vert lagt til side, har ingen stad å stå når det skal seie at \emph{dette} skal ikkje lagast. Det manglar setninga, og det som manglar setninga, endar med å byggje tingen.

\section{Tingas politikk var alltid kjend}

Designfaget kan ikkje hevde at denne innsikta er ny eller utilgjengeleg. Langdon Winner spurde alt i 1980 om artefakter har politikk, og svara ja: bruer bygde for låge til at bussar kjem under, kan halde fattige borte frå ein strand like effektivt som ei lov \autocite{winner1980}. Bruno Latour viste korleis vi delegerer moral og makt til ting — fartshumpen som «politimannen» som aldri tek pause \autocite{latour1992missing}. Peter-Paul Verbeek bygde ut ein heil etikk for korleis ting medierer og formar menneskeleg handling \autocite{verbeek2011}. Heile dette feltet — teknologi- og vitskapsstudia — har i meir enn førti år levert designfaget den innsikta det treng for å forstå at forma er moralsk lada. Innsikta finst. Ho er berre ikkje innbygd i faget sjølv.

\section{Etikken som vart bolta på}

Faget sitt svar er å leggje til eit emne. Det finst no etikk på pensum, berekraft på pensum, klima på pensum; det er pålagt, og det er meint vel. Men sjå kva det inneber at etikken er eit \emph{tillegg}. Det inneber at faget sjølv — modellen av kva design er — er etisk nøytralt, og at moralen kjem utanfrå og vert lagd oppå. Det er feil veg. I kvart fag som har eit fundament foreinleg med korleis verda heng saman, fell etikken ut av faget sjølv: økologien treng ikkje eit påbolta etikkemne for å vite at eit inngrep i éin node forplantar seg gjennom heile nettet; det \emph{er} økologien. Eit kvart fundament foreinleg med samtidas forståing av økologi, biologi og kognisjon ville uunngåeleg produsere ein etikk som la avgrensingar på designaren — ikkje som tillegg, men som konsekvens. At designfaget i staden må bolte etikken på utanfrå, er sjølve beviset på at faget sin modell av seg sjølv ikkje er foreinleg med korleis verda heng saman.

Røystene som har kravd noko anna, kjem stort sett frå randa av faget, ikkje frå kjernen. Mike Monteiro har bedt designarar innsjå at dei har øydelagt mykje og må organisere seg som ein profesjon med ansvar \autocite{monteiro2019}. Sasha Costanza-Chock har lagt fram ein heil ramme for «design justice» som set makt og fordeling i sentrum \autocite{costanzachock2020}. Dette er viktige bidrag, men legg merke til kva dei er: appellar om at faget \emph{bør} ta inn ein etikk det ikkje alt har. Dei stadfestar fråvêret dei prøver å bøte på. Eit fag med eit fundament treng ingen som ber det om å bli etisk; etikken er alt der, i modellen.

\section{Fråvêret av eit organ}

Den endelege forma på dette fråvêret er institusjonell. Eit fag som hadde eit fundament, kunne gjort det individuelle nei om til eit fagleg nei — slik medisinen kan frata ein lege retten til å praktisere, slik ingeniørstanden kan utelukke for grov aktløyse. Designfaget har ingen slik mekanisme, fordi ein slik mekanisme krev ein standard å utelukke for brot på, og faget har ingen standard. Det finst ikkje noko internasjonalt etisk råd for design med makt til å ekskludere. Du kan ikkje miste autorisasjonen i eit fag som ikkje har ein. Difor er dette det tyngste leddet i klagemålet: faget likar å sjå seg sjølv som ei kraft for det gode — menneskesentrert, empatisk, problemløysande — men ein profesjon utan evne til å seie nei er ikkje ei kraft for det gode. Han er ei kraft, ubunden, til disposisjon for den som betalar, og resultatet ber signaturen til kven som helst som leigde han.

\section{Eit barn som eit punkt i ein database}

Lat oss gå heilt ned i det konkrete, for det er der argumentet anten held eller fell. Tenk deg ein analytikar framføre ein skjerm i ein bygning i Washington-området. Grensesnittet ho ser på, er reint, roleg, godt teikna: eit kart med klikkbare nodar, ei venstrekolonne med filter, ei søkjelinje som autofullfører namn. Ho skriv inn ein adresse. Programmet teiknar opp eit nett av samband — kven som bur der, kven dei ringjer, kven dei arbeider med, kva bilar som er registrerte på adressa, kva skular borna går på. Kvar node er eit menneske. Kvar line er eit høve til å hente nokon. Dette er ikkje science fiction; det er produktet Palantir bygde for amerikansk grensekontroll, dokumentert i detalj av journalistar over fleire år \autocite{biddle_intercept}. Systemet heitte fyrst FALCON, seinare Investigative Case Management; det samla data frå eit halvt dusin føderale register i éin søkbar flate, og det vart brukt i operasjonar der born vart registrerte som inngang til å nå foreldra deira — barnet som eit \emph{lead}, ein tråd å dra i.

Sjå no på kvar avgjerda om forma vart teken, og av kven. Nokon valde at autofullføringa skulle vere rask og venleg, for det reduserer friksjon for analytikaren. Nokon valde at sambandsgrafen skulle teikne born og vaksne med same nodetype, for det er reinare i datamodellen. Nokon valde fargane, marghøgda, ikonografien. Kvar einaste av desse er eit formval, teke av ein designar som kunne grunngje det langs aksen «brukarvenleg, effektiv, konsistent» — og som ikkje hadde noko språk i faget sitt for aksen «dette barnet skal ikkje vere ein node». Det er nett det poenget kapitlet held fram: ikkje at designarane var vonde, men at modellen dei arbeidde innanfor, ikkje inneheldt setninga som ville late dei sjå at her var ei grense. Eit reint grensesnitt over ein deportasjonsmaskin er framleis eit godt grensesnitt, etter den einaste målestokken faget eig. Skaden er ikkje at handverket svikta. Skaden er at handverket lukkast.

Og dette er ikkje ein einsleg avsporing. Same arkitektur — eit formrom optimert langs éin akse, navigert av ein agent blind for resten — dukkar opp att der innsatsen er endå høgare. I rapporteringa om målval-systemet «Lavender» kjem det fram korleis ein database skåra titusenvis av personar som potensielle mål, og korleis menneska i lykkja vart reduserte til eit gummistempel som brukte sekund på kvar «godkjenning» før forma slo til \autocite{abraham2024lavender}. Avstanden frå den venlege autofullføringa til dette er kortare enn faget likar å tru. Begge er grensesnitt som gjer ei monstrøs handling rask, lesbar og emosjonelt fjern. Begge er, i streng fagleg forstand, vellukka design.

\section{Augneblinken då ein kunne ha sagt nei}

Det finst ein motsetnad til denne mangelen, og han er verd å sjå like nøye på, for han viser kva eit nei kostar og kva det krev. I 2020 fall det ein dom i Haag som stansa eit system kalla SyRI — Systeem Risico Indicatie — som den nederlandske staten hadde brukt til å skåre innbyggjarar i fattige bydelar for sannsynet for trygdesvindel. Systemet kopla saman opplysningar om bustad, arbeid, gjeld, utdanning og helse, og spytta ut risikoflagg som førte til etterforsking av folk som aldri fekk vite at dei var skåra. Retten kom til at SyRI braut menneskerettane — at det var ein uakseptabel inngrep i privatlivet å la ein ugjennomsiktig modell sortere borgarar i mistenkte og reine — og kravde at systemet vart lagt ned. Det var ikkje ein designar som sa nei. Det var ein domstol. Mekanismen som faget ikkje har innanfor seg sjølv, måtte hentast utanfrå, frå retten, for retten har det faget manglar: ein standard å måle inngrepet mot, og makt til å stanse det som bryt standarden.

Hald desse to bileta opp mot kvarandre. På den eine sida analytikaren framføre Palantir-skjermen, som ikkje ein einaste stad i arbeidsflyten møter ein knapp, ei rute, eit fagleg krav som tvingar fram spørsmålet «skal dette i det heile lagast?». På den andre sida ein dommar som måtte gjere jobben faget ikkje hadde rusta seg til å gjere sjølv. Skilnaden er ikkje moral; eg trur ikkje analytikaren var verre enn dommaren. Skilnaden er at dommaren stod i ein tradisjon med eit fundament — eit korpus av rettar, ein standard, eit organ med makt — og difor kunne forvandle eit privat ubehag til eit bindande nei. Designaren stod i eit fag utan noko av dette, og difor vart ubehaget hennar, om ho kjende det, privat, veglaust, utan adresse å sendast til.

Det er her kapitlet sitt klagemål får si skarpaste form. Ein kunne tenkje at den enkelte designaren berre måtte vere modigare, seie opp, nekte. Men individuell heltemodig nekting er nett det eit fag utan fundament etterlèt mennesket med: ein einsam, ubeskytta, lett utbyteleg gest. Sluttar du, sit den neste i stolen din i morgon, og forma vert bygd likevel. Eit fag med eit fundament gjer noko anna med nei-et: det løftar det frå den private samvitskvalen opp til ein delt standard, slik at det ikkje er du åleine som nektar, men faget som forbyr. Skilnaden mellom SyRI-dommen og Palantir-skjermen er skilnaden mellom eit nei som ber, og eit nei som ingen stad finst.

\begin{kasus}{Den venlege autofullføringa}
Ein interaksjonsdesignar — kall henne M — fortel om eit oppdrag ho fekk tidleg i karrieren, for ein kunde i offentleg sektor ho ikkje fekk namngje. Oppgåva var avgrensa og kjend: betre søkjefunksjonen i eit saksbehandlingsverktøy. Brukarane klaga over at det tok for mange klikk å finne ein person i registeret. M gjorde det ho var lærd å gjere. Ho intervjua brukarane, teikna brukarreiser, kartla friksjonen, og kom fram til at ei autofullførande søkjelinje med samtidig framvising av relaterte oppføringar ville kutte tida frå oppslag til funn dramatisk. Ho prototypa, testa, itererte. Tilbakemeldingane var glitrande: saksbehandlarane elska henne for det, oppgåvetida fall, leiinga skrytte av prosjektet på eit allmøte.

Det var fyrst lenge etter, fortel ho, at ho forstod kva «relaterte oppføringar» eigentleg hadde gjort. Funksjonen hennar synleggjorde slektningar, kontaktar og adressehistorikk for kvar person som vart slått opp — og registeret gjaldt innvandrarsaker. Ho hadde, med dei reinaste handverksmessige instinkt, bygd ein effektiv veg frå éin person til heile nettverket rundt vedkomande. «Eg hadde alle orda for kvifor søket var godt,» seier ho. «Eg hadde ikkje eitt ord for å spørje om søket burde finnast.» Det er heile diagnosen i éi setning. M var ikkje aktlaus; ho var dyktig. Faget hadde gjeve henne eit rikt språk for den eine aksen — effektivitet, brukarglede, redusert friksjon — og ikkje noko språk i det heile for aksen ho seinare ikkje fekk sova for. Ho leita etter ein stad å melde uroa si og fann ingen: inga fagleg norm broten, inga line ho kunne peike på og seie «her, dette forbyr det». Det fanst inga line. Det er det eit fag utan fundament er: eit fag der den dyktige gjer skaden best.
\end{kasus}

\vignett

Korleis kan eit fag i denne tilstanden halde fram utan at samanbrotet vert openbert? Svaret er at fråværet sjølv er gjort usynleg — og at det å gjere fråvær usynleg er det institusjonane i faget eigentleg er gode til. Det neste kapitlet samlar fråværa og ser dei i kvitauget.

\laeringsmal{\begin{itemize}
\item Forklare kvifor mangelen på fundament er eit moralsk problem, og ikkje berre eit epistemisk.
\item Bruke teknologi- og vitskapsstudia (Winner, Latour, Verbeek) til å vise at form er moralsk lada.
\item Gjere greie for kvifor ein etikk bolta på som eige emne er strukturelt feil veg, og kva eit fundament ville gjort i staden.
\end{itemize}}

\ovingar{\begin{enumerate}
\item «Palantir er eit tenestedesignprosjekt.» Forsvar eller motbevis at dette er design i fagleg forstand. Kva står på spel i svaret?
\item Winner spør om artefakter har politikk. Vel eit kvardagsobjekt (ein fartshump, ein sjølvbeteningskasse, eit skjema) og analyser kva politikk det ber.
\item Kvifor kan ikkje designfaget frata ein utøvar retten til å praktisere, slik legestanden kan? Kva nøyaktig måtte vere på plass for at det skulle kunne?
\end{enumerate}}

\vidarelesing{\fullcite{winner1980}; \fullcite{latour1992missing}; \fullcite{verbeek2011}; \fullcite{eubanks2018}; \fullcite{monteiro2019}; \fullcite{costanzachock2020}.}
